\section{Introduction}\label{sec:introduction}

A permutation constraint prescribes exactly one label at one endpoint for
each label at the other. The Unique Games Conjecture asserts that systems
of such constraints can be hard to approximate even when almost all
constraints can be satisfied. We resolve this conjecture positively, with an
explicit unweighted reduction from \(\mathrm{3SAT}\).

\subsection{The theorem}

A \emph{Unique Games instance} has a finite vertex set, a finite alphabet
\(\Sigma\), and a nonempty list of oriented constraints
\(e=(u_e,v_e,\pi_e)\), where \(\pi_e\) is a permutation of \(\Sigma\).
A labeling \(a\) satisfies \(e\) when
\(a(v_e)=\pi_e(a(u_e))\). Its value is the fraction of satisfied
constraints, and \(\val(G)\) is the maximum value of a labeling.
The instance is \emph{explicit and unweighted} when the entire vertex
and constraint lists are written out, including the table of each
permutation. A simple bipartite instance has two disjoint vertex sets
and at most one constraint on each pair of opposite vertices.

\begin{theorem}[Unique Games Theorem]\label{thm:ugc}
For every fixed \(\varepsilon,\delta\in(0,1/2)\), there are an integer
\(s\ge1\) and a deterministic polynomial-time reduction that maps each
\(\mathrm{3SAT}\) formula \(\varphi\) to an explicit unweighted Unique
Games instance \(G_\varphi\) over \(\Sigma=\F_2^s\), with
\[
\begin{aligned}
 \varphi\text{ satisfiable}
 &\quad\Longrightarrow\quad \val(G_\varphi)\ge1-\varepsilon,\\
 \varphi\text{ unsatisfiable}
 &\quad\Longrightarrow\quad \val(G_\varphi)\le\delta.
\end{aligned}
\]
The graph is simple and bipartite, and every constraint is a translation
\(a(v_e)=a(u_e)+c_e\) of \(\Sigma\). The alphabet, the degree of the
running-time polynomial, and its constants depend only on
\(\varepsilon,\delta\).
\end{theorem}

Thus the result resolves the Unique Games Conjecture positively in its
usual edge-value formulation~\cite[Conjecture~2.5]{Khot10}.
The two errors are prescribed independently, before the alphabet is chosen.
Binary translation constraints are a known equivalent formulation of the
conjecture~\cite[Corollary~13 and Section~11.1]{KKMO07}; the reduction
here produces that form directly.

Theorem~\ref{thm:ugc} supplies the Unique Games hypothesis in earlier
hardness reductions for cut, covering, constraint-satisfaction, ordering,
deletion, and clustering problems
\cite{KKMO07,KhotRegev2008,Raghavendra08,GuruswamiEtAl2011OrderingCSP,CKKRS2006,DemaineEtAl2006Clustering}.
The reductions and the resulting approximation thresholds are due to
those authors. Section~\ref{sec:consequences} gives the precise problem
formulations, strict thresholds, and reduction qualifications, together
with proofs of the consequences.

\subsection{Proofs, codes, and the quantifiers}

The central difficulty is to retain completeness near one when every
label permits only one reply. The earlier connection between proof
verification and approximation hardness was developed by Feige,
Goldwasser, Lov\'asz, Safra, and Szegedy~\cite{FGLSS96}, by Arora and
Safra through proof composition~\cite{AroraSafra98}, and by Arora,
Lund, Motwani, Sudan, and Szegedy through the constant-query PCP
theorem~\cite{ALMSS98}. Bellare, Goldreich, and Sudan introduced the
long code and folding for verifier composition~\cite[Section~3.3]{BGS98}.
The long code of an answer \(a\) is the table \(f\mapsto f(a)\),
indexed by Boolean functions on the answer set. In H\aa stad's
parity-hardness proof, Fourier analysis recovers an outer-game strategy
from arbitrary proof tables that make his verifier accept sufficiently
often, without first recovering nearby codewords. His nearly satisfiable
parity gap is our starting hardness input~\cite{Hastad01}.

Khot introduced the Unique Games Conjecture in 2002 and gave an algorithm
based on semidefinite rounding~\cite{Khot02}.
Several algorithmic guarantees explain why its order of quantifiers
matters. For a fixed target error, Trevisan combines semidefinite rounding
with graph decomposition to obtain a polynomial-time
guarantee when the input error shrinks at the reciprocal-logarithmic
rate in the instance size~\cite{Trevisan08}. Charikar, Makarychev, and
Makarychev use randomized semidefinite rounding to obtain expected value
\(1-O(\sqrt{\epsilon\log q})\) from
a \(q\)-label instance of value at least \(1-\epsilon\)~\cite{CMM06};
this guarantee loses accuracy as the alphabet grows, whereas the
conjecture permits a large fixed alphabet after the errors are chosen.
Kolla's spectral enumeration gives a running-time bound in terms of the
dimension of a suitable spectral subspace of the label-extended graph,
whose vertices are variable--label pairs~\cite{Kolla11}. Arora, Barak, and Steurer use
spectral decomposition and enumeration to obtain near-satisfying
solutions for sufficiently small fixed completeness error, with a
guarantee independent of the fixed alphabet size~\cite{ABS10}; their
running-time bound is subexponential, not polynomial, for fixed positive
error.

Integrality gaps concern the strength of a relaxation. Khot and Vishnoi
constructed Unique Games with semidefinite value arbitrarily close to
one and integral value arbitrarily close to zero, even with triangle
inequalities~\cite{KV15}. Barak, Gopalan, H\aa stad, Meka, Raghavendra,
and Steurer introduced the short code, based on Reed--Muller codes, as a
shorter substitute for the long code in applications including alphabet
reduction and quantitative integrality gaps~\cite{BGHMRS15}. These constructions constrain particular
relaxations; a hardness reduction requires a separate argument.
Bafna, Barak, Kothari, Schramm, and Steurer obtain polynomial-time
constant-value guarantees for nearly satisfiable translation games
whose graphs have suitable low-degree sum-of-squares hypercontractivity
certificates, including noisy hypercube and noisy Reed--Muller short-code
graphs~\cite{BBKSS21}. The rank-one matrix shortcode used below is a
different graph.

\subsection{Grassmann expansion and the completeness obstacle}

The closest preceding hardness breakthrough concerns 2-to-1 and
2-to-2 games with imperfect completeness. A 2-to-1 constraint is a
projection with exactly two preimages per right label. A 2-to-2
constraint partitions the labels on each side into pairs and matches
these pairs bijectively, allowing the four combinations within a
matched pair. Khot, Minzer, and Safra's STOC~2017 work proposed a
Grassmann-graph reduction whose nonstandard soundness guarantee was
conditional on a structural hypothesis~\cite{KMS17}. Dinur, Khot,
Kindler, Minzer, and Safra's STOC~2018 work proposed a reduction with
ordinary edge-satisfaction soundness, conditional on an agreement
hypothesis~\cite{DKKMS25}; their expansion work formulated the relevant
structural problem and established partial cases~\cite{DKKMS21Expansion}.
Barak, Kothari, and Steurer connected expansion to the needed agreement
test through a binary matrix representation~\cite{BKS18}. The expansion
theorem of Khot, Minzer, and Safra completed this line in
2018~\cite{KMS23}.

These results established hardness for almost satisfiable 2-to-1 and
2-to-2 games and, as a consequence, a Unique Games gap with completeness
close to one half and arbitrarily small soundness. Obtaining Unique
Games completeness arbitrarily close to one was a separate problem.\footnote{A
separate companion proves hardness for exact 2-to-1 games with perfect
completeness and arbitrarily small fixed soundness, over fixed
alphabets~\cite[Theorem~1.1]{OpenAIPerfectTwoToOne2026}.}

The degree-two matrix shortcode of Barak, Kothari, and
Steurer~\cite{BKS18} is the formulation used here: its vertices are
binary matrices and its steps add rank-one matrices. Their two-reply rule
permits either the old answer or its translate by the update's left factor;
requiring equality alone gives a one-half completeness guarantee. They
discussed higher-degree tests as a possible route to improving this
completeness~\cite[Section~1.6]{BKS18}. The inverse statement we use is
the soundness of their equality test~\cite[Hypothesis~2.1]{BKS18}, which
follows from Grassmann expansion; it appears in Theorem~\ref{thm:shortcode}.
Ellis, Kindler, and Lifshitz later gave a shorter expansion proof using
hypercontractivity on spaces of linear maps~\cite{EKL26}; our external
input remains the Khot--Minzer--Safra theorem. Our proof keeps the
established degree-two inverse theorem and changes the noise using a
nonlinear map. The following comparison explains that choice.

\subsection{The construction and proof}

\paragraph{The nonlinear evaluation.}
For the uniform rank-one test, a nonzero evaluation \(M\mapsto Mz\)
changes by \(a(l^\top z)\) under \(M\mapsto M+a l^\top\). Here \(M\)
is a binary \(\ell\)-by-\(m\) matrix, \(z\in\F_2^m\setminus\{0\}\),
and \(a\in\F_2^\ell\), \(l\in\F_2^m\) are independent and uniform.
This evaluation passes the equality test with probability \((1+2^{-\ell})/2\).
We replace this evaluation by a nonlinear one that is almost unchanged
by a different noise law. The noise must still be detected by sufficiently
many independent linear observations: this is what will let soundness
return to the ordinary shortcode test.

For any prescribed \(p_*>0\), Lemma~\ref{lem:equivariant-noise} fixes a rank threshold
\(r_*\) and, for every sufficiently large \(\ell\), constructs a binary
space \(\mathcal V\) containing \(K=\F_2^\ell\), a map
\(C:\mathcal V\to K\), and a noise law \(\mu\) on \(\mathcal V\) such that
\[
 C(x+h)=C(x)+h\quad(x\in\mathcal V,\ h\in K),\qquad
 \Pr[C(x+a)\ne C(x)]\le p_*,
\]
where in the probability \(x\) is uniform and independent of
\(a\sim\mu\). Every family of linear functionals whose restrictions to
\(K\) have rank at least \(r_*\) detects \(a\) with probability at least
\(1/8\). The threshold \(r_*\) depends only on \(p_*\), leaving \(\ell\)
free to meet the later shortcode and folding requirements. For uniform
\(P\in\Hom(\F_2^m,\mathcal V)\), the new evaluation \(C(Pz)\) changes
to \(C(Pz+a(l^\top z))\) under the rank-one update. Since \(Pz\) is
uniform and \(l^\top z\) is an independent fair bit, its equality test
fails with probability at most \(p_*/2\).

The separation between nonlinear stability and linear detection is proved
in Section~\ref{sec:equivariant-noise}. A quadratic block over a finite field gives
a nonzero output perturbation a fixed probability of disappearing.
Concatenating these blocks, with noise on one random leaf, makes the
nonlinear change probability decay geometrically with the depth
(Lemma~\ref{lem:equivariant-noise-noise}). For linear detection, we track the rank
of a subfamily of observations down the same branch. Random changes of
coordinates make the surviving space of linear functionals likely to satisfy the
following bound: at rank \(r\), restriction loses at most one dimension,
with loss probability at most a constant times \(r\) times the block's
cancellation probability. Taking the field large controls the exceptional
spaces. A one-dimensional loss reduces the sum
\(\sum_{j=1}^r1/j\) by \(1/r\), cancelling that rank factor.
Starting with sufficiently large rank therefore permits enough levels
to make the nonlinear change small while leaving positive rank with
constant probability. A surviving linear functional detects the uniform leaf
noise with constant probability
(Proposition~\ref{prop:equivariant-noise-all-extensions-detection}). Finally, copies
indexed by injections of the smaller output space into \(K\) allow
\(\ell\) to grow with \(r_*\) fixed: a high-rank family restricts onto
the whole smaller dual for most injections
(Section~\ref{subsec:equivariant-noise-enlargement}).

\paragraph{The permutation test and its completeness.}
Section~\ref{sec:game} uses this evaluation to encode tuples
\(U=(e_1,\ldots,e_k)\) of binary parity equations on three distinct
variables. A table is an affine map from the product of their solution
planes to \(\mathcal V\), represented by a matrix with \(m=1+2k\)
columns, including the intercept. We adapt Khot--Safra's virtual
tables~\cite[Section~1.1]{KS13} by identifying tables that induce the
same affine map on the same retained coordinates. Thus a projected
table and its pullback have the same answer. Folding by constant
\(K\)-translations, following Bellare--Goldreich--Sudan~\cite[Section~3.3]{BGS98},
assigns one label to each translation orbit. Checking equality of the
restored answers at \(P\) and \(P+a l^\top\) is then a translation
constraint on those labels. The equivariance of \(C\) makes nonlinear
evaluation compatible with these labels. A global parity assignment of error \(\xi\) supplies
one common evaluation vector on every tuple it satisfies, so the matrix
game has value at least \(1-k\xi-p_*/2\)
(Lemma~\ref{lem:outer-completeness}).

\paragraph{Soundness: reduction to the shortcode test.}
For soundness, suppose a labeling passes the matrix test with probability
at least \(0.99\). A matrix frequency consists of one linear functional
per column. Averaging over \(l\) makes its noise eigenvalue the probability
that all these functionals vanish on \(a\). The detection bound makes
this nonnegative eigenvalue at most \(7/8\) when their restrictions to
\(K\) have rank at least \(r_*\). High acceptance therefore forces a
fixed amount of Fourier mass below that rank. Under uniform noise in
\(K\), the eigenvalue at restriction rank \(r\) is \(2^{-r}\), so this
mass gives positive acceptance of the ordinary rank-one test
(Lemma~\ref{lem:matrix-spectral}). Split \(\mathcal V=K\oplus W\) and
write \(P=(M,T)\) for its two components. The ordinary test fixes \(T\);
on a positive fraction of pairs \((U,T)\), Theorem~\ref{thm:shortcode}
therefore gives agreement with an affine evaluation \(Mz+u\), for some
\(z\in\F_2^m\) and \(u\in K\), on a nonempty slice defined by at most
\(r_{\rm s}\) row equations and \(r_{\rm s}\) column equations, for a
fixed \(r_{\rm s}\).

\paragraph{Row advice and agreement on slices.}
To turn these slices into outer-game strategies, we need a positive
fraction of advice admitting agreement on a slice. The equation prover's advice is
\((A,T,AM)\), where \(A:K\to\F_2^{r_{\rm s}}\) is independently
uniform; it specifies the row fiber \(\{M':AM'=AM\}\).
Lemma~\ref{lem:row-advice-agreement} shows that a fixed positive fraction of
such advice admits a slice with substantial agreement. Its erasure
argument is the key: if the union of row fibers admitting such agreement were too small,
random replacements of their entries could be chosen to preserve enough
shortcode acceptance while destroying the agreement promised by the
inverse theorem on every slice.
The inverse theorem applies to the modified function and gives a
contradiction. Each such advice then determines a fixed choice of
slice and target. Writing its column equations as \(Mq_j=t_j\), we can
vary \(z\) by \(\Span\{q_j\}\) while absorbing the change into \(u\).
Folding forces the resulting coset to contain valid answers to \(U\);
these give a nonempty list of at most \(2^{r_{\rm s}}\) answers.

\paragraph{Transfer to sparsely projected questions.}
We compare \(U\) with a question \(O\) obtained by independently
replacing each equation by one uniformly chosen variable with probability
\(\beta\). This uses Khot--Safra's smooth outer-game method
\cite[Section~3 and Lemma~3.1]{KS13}, with correlated advice as in
Khot--Minzer--Safra~\cite[Sections~3.3--3.4]{KMS17}, also used by
Minzer--Zheng~\cite[Section~3]{MZ26}. The referee samples a uniform
table on \(O\), gives its \(W\)-component and random linear combinations
of its \(K\)-valued rows as advice there, and gives their pullbacks as advice
on \(U\), with the same map \(A\). With all spaces fixed, take \(k\)
large and \(\beta=k^{-2/3}\). Since \(k\beta^2\to0\),
the full pulled-back table law, averaged over projection choices,
approaches the unrestricted uniform law
(Lemma~\ref{lem:projected-tv}). A chosen slice can become empty on
projection; comparing the full table laws transfers both membership in
that slice and agreement with its target. Consequently the chosen slice
survives and retains agreement with positive probability
(Lemma~\ref{lem:projected-witness}).

\paragraph{Decoding.}
On advice admitting such a slice, the \(U\)-strategy samples uniformly from its list. Adapting
Khot--Safra's Fourier-list decoding~\cite[Section~1.1]{KS13}, the
\(O\)-strategy samples Fourier mass on its own row fiber. Agreement on
the transferred slice puts positive mass on frequencies that decode into
the projected list. Each strategy uses only its own question and advice; the
transferred slice certifies their success in the analysis. The
independent sample from the bounded list then gives full answer agreement
with probability at least a constant \(\gamma>0\), independent of \(k\)
(Lemma~\ref{lem:fourier-coset-decoding} and
Proposition~\ref{prop:matrix-decoding}).

\paragraph{The bound from vanishing advice slopes.}
On a NO parity instance, the advice game has vanishing value.
Section~\ref{sec:zero-advice} restricts to singleton coordinates where all
advice slopes are zero. Each singleton has a fixed positive probability
that its advice slopes vanish, and the same sparse choice of \(\beta\) leaves
\(k\beta=k^{1/3}\to\infty\). These coordinates contribute neither
slopes nor intercept corrections to the advice. Fixing \(A\), the positions of these
coordinates, the original intercept, and the other coordinate data therefore
leaves independent equation--variable questions. Each prover can reconstruct
its original input from its own questions on these coordinates and the fixed data, so
parallel repetition applies to these coordinates. Averaging over their
number bounds full advice agreement by \(\exp(-c k^{1/3})\), for a
constant \(c>0\) depending only on the fixed advice spaces
(Lemma~\ref{lem:zero-advice}). This contradicts the decoded lower bound for
large \(k\), giving the fixed matrix-game gap in
Proposition~\ref{prop:matrix-gap}.

\paragraph{Completing the reduction.}
Section~\ref{sec:finish} rounds the weights and replaces each constraint
occurrence by a path of four edges, producing a simple bipartite game
while preserving a fixed gap and the translation form. The final step
uses the same projection-game repetition bound as the argument using vanishing advice
slopes, namely Dinur--Steurer's alphabet-independent bound
\cite[Corollary~1.2]{DS14}.\footnote{Raz proved exponential decay for
general two-prover games~\cite{Raz98}, and Holenstein simplified the
proof~\cite{Holenstein2009}. Those bounds depend on the answer alphabet;
Rao removed that dependence for projection games~\cite[Theorem~4]{Rao11}.}
Its exponent \(t\) can be chosen from the fixed gap and \(\delta\) before
the answer alphabet \(K\); the \(t\)-fold repeated game then has the
alphabet \(\Sigma=K^t\) of Theorem~\ref{thm:ugc}. We then choose the nonlinear change bound and alphabet,
the tuple length \(k\), and finally the parity error, making the
completeness losses small enough for the prescribed \(\varepsilon\).
All dimensions and repetition counts are fixed by the two errors, so
enumerating the finite experiments gives the deterministic polynomial-time
reduction of Theorem~\ref{thm:ugc}.

Figure~\ref{fig:overview} displays the reduction and the two incompatible
bounds that establish soundness.
\begin{figure}[htbp]
\centering
\begin{tikzpicture}[
  x=1cm,y=1cm,
  every node/.style={font=\small},
  box/.style={draw=linkblue!80!black,rounded corners=2pt,
    text width=5.9cm,align=center,inner sep=8pt,minimum height=11mm},
  branch/.style={box,text width=2.70cm,inner sep=6pt,minimum height=15mm},
  heading/.style={font=\small\bfseries,align=center},
  reduction/.style={-{Latex[length=2.1mm]},thick,draw=linkblue!80!black},
  implication/.style={-{Latex[length=2.1mm]},thick,dashed,draw=black!70}
]
\node[heading] at (0,0) {The reduction};
\node[heading] at (7.4,0) {Soundness on a NO input};

\node[box] (parity) at (0,-1.0)
  {$\mathrm{3SAT}$\\$\longrightarrow$ weighted parity gap instance};
\node[box] (matrix) at (0,-2.8)
  {Matrix equality test\\
   stable nonlinear map: Lemma~\ref{lem:equivariant-noise}};
\node[box] (ugc) at (0,-6.2)
  {Unweighted simple bipartite\\Unique Games\\
   Theorem~\ref{thm:ugc}};
\draw[reduction] (parity) -- (matrix);
\draw[reduction] (matrix) -- node[midway,fill=white,inner sep=4pt,
    text width=5.0cm,align=center]
  {Rounding and subdivision;\\parallel repetition} (ugc);

\node[branch] (assumption) at (5.65,-1.0)
  {Assume test acceptance\\at least $0.99$};
\node[branch] (no) at (9.15,-1.0)
  {NO parity\\instance};
\node[branch] (decoding) at (5.65,-3.3)
  {Decoded advice strategies\\agreement at least $\gamma>0$\\
   Proposition~\ref{prop:matrix-decoding}};
\node[branch] (zero-advice) at (9.15,-3.3)
  {Every pair of\\strategies:\\agreement at most\\$\exp(-c k^{1/3})$\\
   Lemma~\ref{lem:zero-advice}};
\node[box] (contradiction) at (7.4,-6.2)
  {Contradiction for sufficiently large $k$:\\
   $\exp(-c k^{1/3})<\gamma$};
\draw[implication] (assumption) -- (decoding);
\draw[implication] (no) -- (zero-advice);
\draw[implication] (decoding) -- (contradiction);
\draw[implication] (zero-advice) -- (contradiction);
\end{tikzpicture}
\caption{The solid arrows describe the reduction; the dashed arrows describe
its soundness argument. The decoding lower bound follows from high test
acceptance, while the upper bound from vanishing advice slopes follows from the NO
parity gap and applies to every advice strategy. The constants
$c,\gamma>0$ are fixed before the tuple length $k$ grows.}
\label{fig:overview}
\end{figure}


\paragraph{Organization.}
Section~\ref{sec:preliminaries} fixes the conventions and the repetition
bound. Section~\ref{sec:game} states the properties of the equivariant map with additive noise and constructs the
permutation test. Section~\ref{sec:matrix-gap} states the two bounds on
advice agreement and deduces the matrix-game gap; Section~\ref{sec:finish}
then fixes the parameters and proves Theorem~\ref{thm:ugc} from these inputs.
The three supporting branches are proved in Sections~\ref{sec:equivariant-noise},
\ref{sec:matrix}, and~\ref{sec:zero-advice}: the construction of the equivariant map with additive noise, matrix
decoding, and the bound from vanishing advice slopes, respectively.
Section~\ref{sec:consequences} gives the approximation consequences.
Appendix~\ref{app:finite-consequences} proves the simple formulation with completeness on left
vertices and the finite rational conversions used there.
